\documentclass[a4paper,12pt]{article}
\usepackage{graphicx, setspace, array, xcolor, mathtools, contour, tikz, amsthm, setspace, amsmath,amsfonts,amssymb, subcaption, fancyhdr, lipsum,multicol, geometry, gensymb}
\usepackage{colortbl}
\geometry{a4paper, margin=0.75in}
\contourlength{0.5pt}
\usetikzlibrary{patterns}
\usepackage[english]{babel}
\renewcommand{\qedsymbol}{$\blacksquare$}
\definecolor{darkgreen}{rgb}{0.0, 0.5, 0.0}
\pagestyle{fancy}
\fancyhf{}
\renewcommand{\headrulewidth}{0.4pt}
\fancyhead[L]{\rightmark}
\fancyhead[R]{\thepage}
\title{Modern Algebra}
\author{Assignment 11\\ \\ Yousef A. Abood\\ \\ ID: 900248250}
\date{December 2025}
\setlength{\parindent}{0pt}
\singlespacing
\parskip=1mm
\setcounter{section}{10}
\setcounter{subsection}{1}
\onehalfspacing
\begin{document}
\maketitle
\noindent\makebox[\linewidth]{\rule{15cm}{0.4pt}}
\section{Finite Abelien Groups}
\section*{Problem 2}
We want to find an integer $n$ such that it is the least integer which we can find three different partitions for the powers of the primes in its prime factorization. To find such an integer, choose the least prime number, which is $2$. Since $3$ has three partitions, we have three options:
\begin{itemize}
    \item partition 3: $\mathbb{Z}_{2^3}$
    \item partition 2+1: $\mathbb{Z}_{2^2}\oplus \mathbb{Z}_2$
    \item partition 1+1+1: $\mathbb{Z}_2\oplus \mathbb{Z}_2\oplus \mathbb{Z}_2$
\end{itemize}
\section*{Problem 5}
By the Fundamental Theorem of Finite Abelian Groups, observe that an abelian group of order $45$ is isomorphic to $G_1\cong \mathbb{Z}_5 \oplus \mathbb{Z}_{3^2}$ or $ G_2 \cong \mathbb{Z}_5\oplus \mathbb{Z}_3 \oplus \mathbb{Z}_3.$ We know that the maximum order of an element in $G_1,G_2$ is $\operatorname{lcm}(9,5)=45, \operatorname{lcm}(3,5)=15$ respectively. Since $15$ divides both $15,45$. Thus, the first statement is true. For the second statement, since $9 \nmid 15$, not all abelian groups of order $45$ has an element of order $9.$
\section*{Problem 10}
We know that the prime factorization of $360$ is $5\cdot 3^2 \cdot 2^3.$ By the Fundamental Theorem of Finite Abelian Groups, any abelian group of order $360$ is isomorphic to one of the following partitions:
\begin{itemize}
    \item $\mathbb{Z}_{5} \oplus \mathbb{Z}_{3^2} \oplus \mathbb{Z}_{2^3}$
    \item $\mathbb{Z}_{5} \oplus \mathbb{Z}_{3} \oplus \mathbb{Z}_3 \oplus \mathbb{Z}_{2^3}$
    \item $\mathbb{Z}_{5} \oplus \mathbb{Z}_{3^2} \oplus \mathbb{Z}_{2^2} \oplus \mathbb{Z}_{2}$
    \item $\mathbb{Z}_{5} \oplus \mathbb{Z}_{3^2} \oplus \mathbb{Z}_{2} \oplus \mathbb{Z}_{2} \oplus \mathbb{Z}_{2}$
    \item $\mathbb{Z}_{5} \oplus \mathbb{Z}_{3} \oplus \mathbb{Z}_{3} \oplus \mathbb{Z}_{2^2} \oplus \mathbb{Z}_{2}$
    \item $\mathbb{Z}_{5} \oplus \mathbb{Z}_{3} \oplus \mathbb{Z}_{3} \oplus \mathbb{Z}_{2} \oplus \mathbb{Z}_{2} \oplus \mathbb{Z}_2$
    

\end{itemize}
\section*{Problem 22}
\begin{proof}
    By the Fundamental Theorem of Finite Abelian Groups, we have a sequence of primes $p_1, p_2, p_3, \cdots , p_k$ and a sequence of integers $n_1, n_2, n_3, \cdots , n_k$ such that
    \[G \cong \mathbb{Z}_{p_1^{n_1}} \oplus \mathbb{Z}_{p_2^{n_2}} \oplus \mathbb{Z}_{p_3^{n_3}} \oplus \cdots \oplus \mathbb{Z}_{p_k^{n_k}}\]
    To prove that $G$ is cyclic, we must show that the primes $p_1, p_2, p_3, \cdots , p_k$ are all distinct. 
    Assume for the sake of contradiction that the group is not cyclic, so we have $p_1=p_2=p.$
    We find $H_1=\mathbb{Z}_{p^{n_1}}\oplus \{0\} \oplus \{0\}\oplus \cdots \oplus \{0\}, H_2= \{0\}\oplus \mathbb{Z}_{p^{n_2}} \oplus \{0\}\oplus \cdots \oplus \{0\}$ are both cyclic subgroups in $G.$ Clearly, each of $H_1, H_2$ has a subgroup of order $p$. Thus, we found two distinct subgroups of order $p$ contradicting that $G$ has exactly one subgroup for each divisor of $|G|.$ Hence, $G$ is cyclic.
\end{proof}
\section*{Problem 31}
By the Fundamental Theorem of Finite Abelian Groups, we know that $G$ is isomorphic to one of five following groups:
\begin{enumerate}
    \item $\mathbb{Z}_{2^4}$
    \item $\mathbb{Z}_{2^3}\oplus \mathbb{Z}_{2}$
    \item $\mathbb{Z}_{2^2}\oplus \mathbb{Z}_{2^2}$
    \item $\mathbb{Z}_{2^2}\oplus \mathbb{Z}_{2}\oplus \mathbb{Z}_{2}$
    \item $\mathbb{Z}_{2}\oplus \mathbb{Z}_{2}\oplus\mathbb{Z}_{2}\oplus \mathbb{Z}_{2}$
\end{enumerate}
Clearly, we eliminate the fifth group as it does not contain elements of order $4.$
We eliminate groups 1, 2, and 4 because in these cases, the squares of all elements of order 4 are identical:
\begin{itemize}
    \item In Group 1 ($\mathbb{Z}_{16}$), the elements of order $4$ are $4$ and $12$. Their squares are $2(4)=8$ and $2(12)=24 \equiv 8$. Thus $a^2 = b^2$.
    \item In Group 2 ($\mathbb{Z}_{8}\oplus \mathbb{Z}_{2}$), any element $g=(x,y)$ of order $4$ must have $x \in \{2,6\}$. The square is $2g = (2x, 2y) = (4, 0)$. Thus $a^2 = b^2$.
    \item In Group 4 ($\mathbb{Z}_{4}\oplus \mathbb{Z}_{2}\oplus \mathbb{Z}_{2}$), any element $g=(x,y,z)$ of order $4$ must have $x \in \{1,3\}$. The square is $2g = (2x, 2y, 2z) = (2, 0, 0)$. Thus $a^2 = b^2$.
\end{itemize}
Observe the third case $\mathbb{Z}_{2^2}\oplus \mathbb{Z}_{2^2}$, consider $a=(1,0), b=(0,1)$ elements of order $4$ in $\mathbb{Z}_{2^2}\oplus \mathbb{Z}_{2^2}$. See that $(0,1)+(0,1)=(0,2)$ and $(1,0)+(1,0)=(2,0)$. Thus, $a^2 \ne b^2$. Hence, $G \cong \mathbb{Z}_{2^2}\oplus \mathbb{Z}_{2^2}.$
\end{document}